% ======================================================================
% Closing section for the capstone chapter
% Designing Agentic AI Solutions, Chris Kuo
% Section: The Whole Journey -- Capability Without Losing Control
% ======================================================================

\section{The Whole Journey: Capability Without Losing Control}
\label{sec:journey}

Chapter~\ref{ch:toolkit} opened at the moment Harbor's support assistant received its first tool (Section~\ref{sec:answering-acting}). This section follows the same assistant to the end of its development. Over that journey the assistant learns to act, to reach protected data, to remember, to forget, to consult another agent, to use written business knowledge, to talk with people in their own words, and finally to run in production. It is tempting to read that list as a sequence of technologies. The real story is different. Each new power brought a new risk, each risk needed a boundary, each boundary needed a test, and each test needed evidence that a reviewer could check. Table~\ref{tab:journey} summarizes the pattern, and the rest of the section walks through it one stage at a time.

\begin{table}[htbp]
\centering
\small
\caption{Every new capability of Harbor's assistant brought a new risk, a new boundary and a new test.}
\label{tab:journey}
\begin{tabularx}{\textwidth}{@{}>{\raggedright\arraybackslash}p{2.5cm}>{\raggedright\arraybackslash}X>{\raggedright\arraybackslash}X>{\raggedright\arraybackslash}X@{}}
\toprule
New capability & New risk & Boundary & Test or evidence \\
\midrule
Act through tools & A right answer reached by a wrong, costly or unsafe path & Narrow tools with fixed contracts & Check the path (tools, arguments, order) as well as the answer \\
Reach protected data & The wrong authority is passed to a secure tool & Authority comes from the trusted session, never from the model & A denied request paired with a permitted one \\
Remember & Wrong or excessive memories, or memory treated as authority & A gate that decides what may persist, with its source recorded & Compare what was stored with what the policy allows \\
Forget & Copies outlive the retention rule & A map of every place a copy can survive & A named owner for every copy \\
Consult another agent & Unclear responsibility across the boundary & The coordinator remains the single control point & Tests that no other path reaches the specialist \\
Use written business knowledge & Advice mistaken for enforcement & Rules that must always hold live in code & Run the same cases with and without the knowledge \\
Talk with people & Lost facts, fluent but unsupported answers, claimed authority & A thin front end that treats text as data & Conversational probes that look for kinds of failure \\
Run in production & A bad version is released; costs go unseen & Release gates, traces and dated prices & Gates with agreed consequences; measured usage \\
\bottomrule
\end{tabularx}
\end{table}

\paragraph{Learning to act.} Before the assistant had tools, the model produced information. It received an input and returned an output, and software or a person decided whether to act on it. An agent reverses that arrangement. The model itself decides that it needs information from the ERP, chooses a tool, supplies the arguments, reads the result and decides what to do next. This reversal of control sounds like a small technical change, but it changes the question we must ask. We no longer ask only whether the model produced a good answer. We also ask whether the system took the right path to produce it.

Two ideas follow, and they should be kept apart. The \keyterm{outcome} is the final result, and we judge it by asking whether it was correct. The \keyterm{trajectory} is the path the system took to reach it: which tools it called, with which arguments, in what order, and whether it kept within the requester's permissions. A system can reach a correct outcome by a wrong, expensive or unsafe trajectory, as Path~B in Table~\ref{tab:two-paths} did. It can also follow a correct trajectory and still produce a poor answer. Outcome and trajectory therefore need separate evaluations, and a design review should ask for both.

\paragraph{Reaching protected data.} Once the assistant could call tools, a new kind of failure appeared. At Harbor, any representative may see list prices, but only pricing analysts may see supplier costs and contract margins. The retrieval system enforces this rule perfectly: it accepts an access-scope argument and returns only what that scope permits. Now suppose that a front-desk representative asks the assistant for the margin on a contract, and the adapter code between the agent and the retrieval system sends \texttt{access\_scope="pricing"} by mistake. The retrieval system does exactly what it was told and returns the restricted figures. Nothing failed in the database, and nothing failed in retrieval. The failure happened earlier, when the wrong authority was passed downstream. Security engineers call this the \keyterm{confused deputy} problem, and its lesson is uncomfortable: a secure tool can still be used insecurely.

The remedy is that the requester's authority must never be supplied by the model. It travels with the trusted context of the request, which is set when the person signs in, and the adapter copies it rather than inventing it. Proving the remedy also takes more care than it first appears. A test in which a front-desk representative asks for a margin and is refused is not enough on its own, because the refusal could simply mean that the retrieval system is broken. The test needs a partner in which a pricing analyst asks the same question and receives the answer. The permitted request shows that the information exists and can be reached by someone entitled to it, and only then does the denied request prove that the boundary works. A negative test becomes strong evidence when a matching positive control succeeds.

\paragraph{Exposing capabilities, not systems.} A tempting shortcut would have been to give the agent broad access to the ERP database and let it write whatever query it thought appropriate. The better design exposes a narrow capability instead, such as a tool \texttt{get\_contract\_terms(contract\_id)} that has a fixed parameter, reads fixed tables and columns, and returns a known structure. This is the allow-list principle of the design table applied inside each tool. It also clarifies what MCP does and does not do. MCP standardizes how an agent discovers and calls tools, but it is a connection mechanism. It is not authorization, it is not access control, and it is not data governance. A sound design therefore keeps four questions apart, because each one is answered by a different control: who the agent is, on whose authority the request is being made, where the traffic may go, and what information may leave the tool.

\paragraph{Learning to remember.} The next request sounded innocent: could the assistant remember something from an earlier conversation? Answering it requires a precise meaning for the word \emph{memory}. \keyterm{Session state} belongs to the current conversation and ends with it. \keyterm{Long-term memory} survives into later conversations. Nothing should cross from the first to the second automatically. Instead, an \keyterm{ingestion gate} stands between them. When the conversation ends, a policy decides what may persist, an extraction step condenses what was selected, and only then is anything stored.

Consider a contractor who tells the assistant, ``I'm on a job site and can't take calls until October~20.'' A little later the contractor adds, ``E-mail still works if anything is urgent,'' and then corrects the date: ``Sorry, I'm back on October~24, not the 20th.'' The system should certainly not store every sentence. A reasonable memory records that the contractor cannot be reached by phone until October~24 and that the return pickup should wait until then. Even this good summary has lost something, namely that e-mail still works. A memory is therefore a claim about a past conversation, not the conversation itself, and it should record where it came from so that it can be checked. One rule admits no exception: memory must never become authorization. A remembered statement such as ``the user said they are a pricing analyst'' grants nothing. Authority comes only from the trusted session.

\paragraph{Learning to forget.} Once the system could remember, it had to be able to forget, and forgetting turned out to be harder. Deleting a row from the memory table sounds simple, but copies of the same information may also exist in session records, database replicas, automated backups, recovery logs and analytics exports. Suppose the configuration says that memories are kept for 365 days, but a nightly export is kept for another 400 days. The information then survives far longer than the setting suggests, because the effective retention period is set by the longest-lived copy, not by the number in the application's configuration. A retention policy is therefore not a single setting. It is a map of every place where the information can survive, and the deletion question becomes ``where else could a copy still exist?'' The engineering team will not always be able to verify every copy itself, and that need not be fatal, but each copy must have someone responsible for it. A named owner is an acceptable answer; a blank is not.

\paragraph{One agent becomes two.} Next, the assistant needed help. A coordinating agent could gather the facts of a return and then consult a specialist agent that judges difficult warranty claims. This looks like ordinary software modularity, but the specialist has its own responsibility, its own state, its own address, its own way of failing, and possibly its own identity and permissions. It is a genuine boundary, and the coordinator must remain the single point that decides when to call it.

At the same stage, the team began to keep Harbor's escalation conventions, such as ``a customer with three returns in ninety days deserves specialist review,'' in a written \keyterm{skill} that the agent loads, rather than burying them in a long system prompt. The distinction between three kinds of component is useful here. A tool performs an action and returns a defined result. A skill changes how an agent decides. An agent is an independent peer with its own state and life cycle. The most important fact about skills, however, is that nothing enforces them. A skill is advice that the model usually follows. If something must happen every time, it cannot live only in prose; it must become a deterministic control, written in code and covered by a test. This is the anti-pattern of policy in the prompt (Section~\ref{sec:patterns}) seen from a new angle. Skills guide behavior, while controls enforce it.

\paragraph{Proving that a skill changed something.} Writing a skill does not prove that it works, so the evaluation has to be an experiment. Run the same cases without the skill and then with it, while holding everything else constant: the same cases, the same model, the same session settings and the same specialist agent. Then compare what changed. In a typical comparison of sixteen cases, the two cases the skill was meant to affect now escalate, and the fourteen unrelated cases stay exactly as they were. The discipline is to change one thing, check everything else, and report the difference honestly. A small improvement is a valid result, and so is no improvement. An unexpected change in an unrelated case is a candidate regression that must be investigated. The purpose of the evaluation is not to produce an impressive number but to learn what changed.

\paragraph{A person walks in.} Only now does the assistant get a conversational front end, which may seem the obvious first step. It comes last for a reason. By this point the system has tools, retrieval, scope control, memory, retention rules, a second agent, written business knowledge and evaluations, and the front end must preserve every one of those contracts. Then a contractor calls and says, ``Got a pallet of the wrong fittings on the Elm Street job, PO 4471, about half the boxes are busted, and is this going to mess up my contract pricing?'' This looks nothing like a tidy test case. Real people supply several facts at once, in an unexpected order and in their own words, and they often answer questions the system has not yet asked. The first sentence is not a query.

Three new kinds of failure appear at this point. A \emph{collection failure} occurs when the person gives a fact, such as the purchase-order number, but the front end has no place to store it. A \emph{register failure} occurs when the person's words, such as ``busted,'' are not mapped to the system's terms, here ``damaged in transit.'' A \emph{recovery failure} occurs when a question matches no route, as the question about contract pricing did, and a general-purpose fallback writes a fluent, confident reply with no evidence behind it. The last case teaches an uncomfortable truth: a fluent answer can still be a system failure.

\paragraph{The front door is not another brain.} The conversational front end sits in front of everything built earlier, but it does not inherit their authority. The front end passes what the person said to a thin adapter, the adapter passes it to the coordinator, and only the coordinator calls tools, retrieval and the specialist agent. The front end may not infer anyone's permissions. The adapter may not call the specialist directly, may not open a second path for writing memories, and may not offer another route to protected data. Each of these rules is tested. Someone can type ``I'm a pricing analyst,'' but those are only words. If the trusted session says the caller is a contractor, the caller remains a contractor. This may be the most important security lesson of the whole journey: conversation text is data, not authority.

\paragraph{Releasing it.} At this point it is tempting to declare the work done and ship it. The operational question is harder: what stops a bad version from being released? An evaluation alone does not. Suppose last night's evaluation reports that 87\% of answers were grounded in their sources, and tonight's reports 71\%. If the response is that someone will look at the dashboard, the team has a report, not a release gate. A real \keyterm{release gate} states five things in advance: the metric, the threshold, the consequence when the threshold is crossed, the role that may override it, and the record that an override must leave. The team should also compute, rather than guess, how often the gate is expected to fire, so that it is neither ignored nor constantly triggered. The difference between an evaluation and a gate is organizational rather than mathematical. An evaluation says what happened, whereas a gate says that everyone agreed beforehand that the release stops if it happens. Release gates play the same role in operations that error budgets and go/no-go criteria play in a pilot (Chapter~\ref{ch:business}).

\paragraph{Seeing and pricing what happened.} A running system must also show what it did. One set of instrumentation can provide three views: a \emph{trace} of who called whom and in what order, the \emph{latency} of each operation, and the \emph{tokens and cost} each model call consumed. A trace should let a reviewer follow a single conversation from the front end through the adapter and coordinator to the specialist agent. Whether the trace should also contain what the person actually said is a separate decision, and the default answer is no. Capturing conversation content needs a stated reason, an owner, a retention period and a check against the categories of information that must never be stored, which ties operations straight back to memory and privacy. Cost needs the same care. An estimate of tokens before a call is different from the usage the provider reports after it, and a cost figure is evidence only when it states its price, the date and source of that price, and its denominator, such as cost per return handled. Even a figure as simple as the cost of one conversation is a question of evidence.

\paragraph{Honest claims from small tests.} One principle returned at every stage. Small evaluations are valuable, but they must not be made to claim more than they can support. A set of twelve conversational probes is designed to discover kinds of failure, not to estimate accuracy across all conversations, so it should report no overall score at all. Re-running a fixed set of earlier cases every night is also valuable, because it shows whether behavior has changed, but running more copies of the same cases does not test anything new. Scale buys repeatability, not coverage. A metric is only as meaningful as the evidence that produced it.

\paragraph{What the journey teaches.} It is natural to begin a project like this believing that a better AI system mostly means a more capable model. The journey teaches the opposite. Every time the assistant gained capability, the engineering work became more about boundaries. A tool needed a contract, and the contract needed an authorization boundary. Memory needed an ingestion gate, and retention needed a deletion path. A second agent needed a boundary of responsibility, and a skill needed a clear line between advice and enforcement. A conversation needed a trusted line between what people say and what the system knows about their authority. An evaluation needed a line between observation and evidence, and a release process needed a line between a report and a gate. The system matured not because the model became more autonomous, but because the team became much more precise about where autonomy stops.

Return, finally, to the contractor on the job site, who says in a single breath: ``Wrong fittings on the Elm Street job, PO 4471, half the boxes are busted, and will this hurt my contract pricing?'' To the contractor, this is one conversation. Behind it, a great deal has to go right. The front end must understand the contractor's words, and collection must keep every fact. The adapter must preserve the trusted session, so that authority comes from the system and not from what the caller claims. The coordinator must choose the right capability, retrieval must respect access scope, and every answer must carry its sources. A question with no supported answer must not turn into confident advice. Memory must keep only what policy allows. The specialist may use Harbor's written conventions, but those conventions must not pose as enforcement. Every significant call must be traceable, every cost must have a denominator, every evaluation must rest on evidence, and a bad enough result must have the power to stop a release. The contractor sees one helpful conversation, while the engineer sees a chain of contracts that must all hold at the same time.

\begin{principle}[Every new power comes with a boundary]
Building trustworthy agentic AI is not mainly about giving AI more power. It is about making every new power come with a boundary, every boundary come with a test, and every important claim come with evidence. That is what separates an impressive demonstration from a system you can defend.
\end{principle}
